\subsection*{Choosing a phrase from the pictured context}
A writer can favour the correct modern phrase when given two alternatives and still produce an inaccurate sentence when asked to translate the passage. We compared positive rehearsal with a loss that also penalised the alternative belonging to the other source (Fig.~\ref{fig:context-phrases}). Both writers improved from 13 to 39 correct training pairs out of 39. A source-ink overlap rule already solved 32. On fifteen new pairs, both trained writers chose both alternatives correctly only for the same change in belief polarity. Neither resolved a new archival pair in both directions. With the complete answer history withheld, each writer began five of the 78 training outputs with the entire supervised prefix and completion.

Both writers scored zero faithful translations on eighteen designated archive readings. The phrase composer scored one, the text teacher eleven and the archive gloss fifteen. The crossed writer conveyed both specified details on three of fifteen pairs. Copying preserved twelve. In one authored clause the crossed writer conveyed that the teacher learns ritual from a pupil. Reversing who asked a question in a father--son exchange produced an answer that failed to identify the son as the respondent. The crossed loss improved neither of the two registered transfer measures. The conditional training success establishes that the pictured source can affect the learned phrase ranking on familiar pairs. It gives little guidance about which relation a new passage expresses.
